\uuid{KhiP}
\titre{Dénombrement}
\auteur{bodin ; hanani ; mzari}
\organisation{exo7}
\licence{CC BY-NC-SA 4.0}
\datecreate{2020-01-20}
\niveau{L1}
\module{Algèbre}
\chapitre{Logique, ensemble, raisonnement}
\sousChapitre{Ensemble}
\difficulte{3}
\theme{Ensembles, applications, dénombrement}
\exo7qcmid{lille-1.ensembles.028}

\contenu{
\begin{question}
\qclassification{100.02,101.01,102.01,102.02}
\qtags{motcle=dénombrement}
On note $\mathrm{C}^k_n$ le nombre de choix de $k$ éléments parmi $n$. Combien fait $\displaystyle \sum _{k=0}^{100}(-1)^k2^{-k}\mathrm{C}^k_{100}$ ?
\begin{answers}
    \bad{$0$}
    \good{$2^{-100}$}
    \bad{$2^{100}$}
    \bad{$100$}
\end{answers}
\begin{explanations}
Utiliser le binôme de Newton, $\displaystyle \sum _{k=0}^{100}\mathrm{C}^k_{100}\left(-\frac{1}{2}\right)^k=\left(1-\frac{1}{2}\right)^{100}=\frac{1}{2^{100}}$.
\end{explanations}
\end{question}
}
